\section{Attribution and verification}
The incidence constructions, target classifications and named
classical designs originate in [L] and its cited predecessors. The
common residue theorem and chamber bijection organize the deductions;
they do not replace the outside classifications of [BDL] and [K]. The
cyclic design used above is classical. The source's exceptional-value,
path-length and numerical corrections were established and independently
reviewed in earlier project work. Their inclusion here supplies one
complete edition, without claiming new historical priority or another
solution of the unital conjecture.

The finite calculations use the actual native Mathieu actions. Their
role is to verify design blocks, complete automorphism groups, explicit
derived-plane maps and stabilizer kernels and complements, not to infer
an abstract group or a geometry from its order. The full programs are
included separately. In particular, the full automorphism calculation
uses GAP together with DESIGN and GRAPE; the independently written
finite-field validator checks the saved incidence and action witnesses
but does not itself prove completeness of the automorphism search.
The projective and affine all-parameter proofs, the common residue
recursion, the classification deductions and the Hermitian diameter
proof do not depend on the successful execution of small examples.

The finite Lean certificates verify the three displayed derived
$PG(2,4)$ incidence maps, the literal disconnected residue over
$\mathbb F_4$ and the cyclic biplane counterexample. The Hermitian
certificate verifies the actual plane transport formula and specified
unitary moves; the elementary connectivity theorem is a separate
formal component. The identification of the finite Witt tables with
native group orbits, full automorphism computations, classical
classification inputs and the all-rank geometric deduction are not
formalized by these certificates. Each declaration and its exact scope
are recorded in the accompanying statement map.

The object descriptions and different type-size multisets of the two
rank-five Mathieu residues are proved above, as are the actual
stabilizer groups. Their identification with the earlier geometries
named $\Gamma_{23}(4)$ and $\Gamma_8(4)$ is not proved here. The accepted
principal result also does not extend the source's necessary
classification to arbitrary rank-two incidence structures without
partial linearity; the cyclic counterexample explains why that
qualification matters.

\section{References}
[L] D. Leemans, \emph{A generalization of a construction due to Van
Nypelseer}, Beitr\"age zur Algebra und Geometrie 46 (2005), 561--574.
The fixed journal edition supplies Constructions 1--3 and Theorems
1--11, including the intermediate statements corrected above.

[G] \emph{Flag relations and derived incidence geometries}, the common
paper of this topic. Its Part II proves Theorems S.1--S.6 and Part III
proves Theorems H.1--H.3, with full proofs.

[BDL] F. Buekenhout, P.-O. Dehaye and D. Leemans, \emph{RWPRI and
$(2T)_1$ flag-transitive linear spaces}, Beitr\"age zur Algebra und
Geometrie 44 (2003), 25--46, Theorem 6.2 on pp.42--43.

[K] W. M. Kantor, \emph{Homogeneous designs and geometric lattices},
Journal of Combinatorial Theory, Series A 38 (1985), 66--74,
Theorem 3 and its proof on p.73.

[P] A. Pasini, \emph{Flag-transitive extensions of dual affine planes},
European Journal of Combinatorics 17 (1996), 657--671. The source cites
this work for the two historical residue names. Those names are recorded
as unproved identifications here, not used as premises.
